The system was built incrementally in phases. Each phase had a stated goal, a small
experiment, an artifact saved to disk and a verification gate (Section~\ref{sec:gates}). This
chapter describes the phases in order, then the engineering problems that shaped the
implementation.

\section{Methodology}
The working rule was \emph{baseline first}: no RACS component was wired into the pipeline until
the reproduced baseline passed its gate, so the comparison in Chapter~6 is against a measured
baseline rather than against numbers copied from the paper. Every phase followed the same loop:
state what is built and why; build the smallest experiment; test it; explain the output;
integrate; save the artifact with its configuration hash; move on.

\begin{figure}[H]
	\centering
	\begin{tikzpicture}[scale=0.85, transform shape]
		\hyphenpenalty=10000\exhyphenpenalty=10000
		% baseline milestones, left to right
		\node[blk, text width=2.3cm] (m1) at (0,0) {\textbf{Phases 1--3}\\environment, VAE, diffusion cost};
		\node[blk, text width=2.3cm] (m2) at (3.0,0) {\textbf{Phases 4--6}\\corpus, CLIP, frozen splits};
		\node[blk, text width=2.3cm] (m3) at (6.0,0) {\textbf{Phases 7--9}\\$K_{\mathrm{pub}}$, $X_{\mathrm{feat}}$, references};
		\node[blk, text width=2.3cm] (m4) at (9.0,0) {\textbf{Phases 10--14}\\EDICT, hiding, channel, CodeFormer};
		\node[blk, text width=2.3cm] (m5) at (12.0,0) {\textbf{Phases 15--16}\\baseline frozen: quality, security};
		% RACS milestones, right to left underneath
		\node[newblk, text width=2.3cm] (m6) at (12.0,-3.0) {\textbf{Phases 17--20}\\retrieval, token pick, rerank};
		\node[newblk, text width=2.3cm] (m7) at (9.0,-3.0) {\textbf{Phases 21--22}\\RACS end to end, comparison};
		\node[newblk, text width=2.3cm] (m8) at (6.0,-3.0) {\textbf{Exact layer}\\lossless, sealed, adaptive embedding};
		\node[src, text width=2.3cm] (m9) at (3.0,-3.0) {Phone prototype};
		\draw[arr] (m1) -- (m2);
		\draw[arr] (m2) -- (m3);
		\draw[arr] (m3) -- (m4);
		\draw[arr] (m4) -- (m5);
		\draw[arr] (m5) -- node[left, font=\scriptsize, align=right]{baseline gate\\passed} (m6);
		\draw[arr] (m6) -- (m7);
		\draw[arr] (m7) -- (m8);
		\draw[arr] (m8) -- (m9);
		\node[font=\footnotesize\bfseries, text=gray!50!black, anchor=south west] at ([yshift=0.15cm]m1.north west) {Baseline reproduction (AgentSemSteCom)};
		\node[font=\footnotesize\bfseries, text=green!35!black, anchor=north west] at ([yshift=-0.15cm]m9.south west) {RACS};
	\end{tikzpicture}
	\caption{Order of work. Every phase ends with an artifact and an automatic verification gate; nothing RACS-specific was wired in before the baseline gate passed.}
	\label{fig:phases}
\end{figure}

\section{Foundation (Phases 1--3)}
\subsection{Phase 1: environment and tensor contract}
A virtual environment with PyTorch~2.6 and CUDA~12.4 was created, and all model caches were
redirected to the data drive. A seeding utility fixes Python, NumPy and PyTorch (CPU and CUDA)
random states. The tensor contract was fixed: images are RGB, \texttt{float32}, shape
$(1,3,H,W)$, values in $[0,1]$ for metrics and $[-1,1]$ for the VAE. The gate checks CUDA
availability, configuration loading, tensor shapes and that two seeded runs produce identical
tensors.

\subsection{Phase 2: VAE and latent space}
The Stable Diffusion VAE compresses a $512\times512\times3$ image to a $64\times64\times4$ latent
(an $8\times$ spatial reduction). A round trip over five test images gave a mean PSNR of
32.81\,dB (range 31.78--34.15) and SSIM 0.935. This sets a ceiling: any pipeline that passes
through the VAE cannot recover a secret better than about 33\,dB.


\subsection{Phase 3: diffusion generation cost}
Stable Diffusion v1.5 was loaded in half precision with attention slicing and a DDIM scheduler.
A sweep over steps and seeds measured cost: at 20 steps one $512$\,px image took 5.56\,s
(0.278\,s per step) with 2.60\,GB peak GPU memory, so the paper's 50 steps cost 13.5\,s per
image. Disabling classifier-free guidance reduced cost by a factor of 1.73, confirming that the
UNet forward pass dominates. Cost grows linearly with the number of denoising steps (5.56,
8.18 and 13.47\,s at 20, 30 and 50 steps), which is exactly the per-message cost that RACS removes.

\section{Retrieval Inputs and Test Harness (Phases 4--6)}
\subsection{Phase 4: retrieval corpus}
A builder deterministically sampled 1{,}000 images from COCO val2017 with a fixed seed,
centre-cropped and resized them to $512\times512$ RGB, and wrote a versioned index file listing
each image identifier, path and SHA-256 hash. The row order of this file is part of the protocol,
because the index returns row numbers. The verification gate checks every hash, the row count
and that a rebuild is identical.


\subsection{Phase 5: CLIP embeddings}
CLIP ViT-B/32 embeds each corpus image to a 512-dimensional, L2-normalised \texttt{float32}
vector. Embeddings are kept in single precision because half precision changed nearest-neighbour
orderings in borderline cases. Re-encoding the corpus gave a maximum absolute difference of
$2.5\times10^{-7}$. Using COCO captions as queries, the correct image was the top-1 result 47.6\,\%
of the time out of 1{,}000; mean similarity was 0.779 within a topic and 0.542 across topics.


An important observation was that CLIP similarity and structural similarity disagree: two
images can be semantically close but laid out completely differently (in one case, CLIP
similarity 0.643 but edge overlap only 0.199). This motivated the structural reranking of Phase 20.


\subsection{Phase 6: evaluation harness and frozen splits}
Two frozen, role-tagged splits were built (Section~3.7). The calibration split (UniStega, 22 per
class) may only be used to check the paper's regime and to tune parameters; the evaluation split
(COCO, Commons portraits, Commons cats, 22 per class, disjoint from the corpus) is the only
source of claims. The split module refuses to return evaluation data to a tuning routine.
Each image is assigned a fixed token stored in a versioned token file. The facial evaluation
images were filtered by title and by CLIP to keep only real single-face photographs. Metrics
(MSE, PSNR, SSIM, LPIPS with an AlexNet backbone) and a frozen report schema were implemented, and
a self-test with 83 checks verifies them.

\section{Baseline Reproduction (Phases 7--16)}
\subsection{Phase 7: public conditions}
For each of the 132 split images, BLIP produces a perception caption with greedy decoding, and
Qwen2.5-0.5B-Instruct rewrites it into a concise generation prompt using a fixed one-shot chat
template. The reply is sanitised (instruction echoes, markdown and URLs removed); a degenerate
reply falls back deterministically to the perception caption, and the number of fallbacks is
recorded (zero occurred). OneFormer produces 150-class ADE20K segmentation for common images.
OpenPose produces body and face keypoints for portraits; on 3 of the 22 tightly cropped
evaluation portraits OpenPose detects no body, and a MediaPipe face mesh is used instead, with
the method recorded per image. To fit in memory, the build runs in two passes (all prompts, then
all features) with one model resident at a time.

\subsection{Phase 8: reference generation}
PicX-Real with the Stable Diffusion v1.5 segmentation and OpenPose ControlNets generates one
reference per image at $256$\,px with 20 DDIM steps and guidance 7.5, from a token-derived seed.
Two independent runs produced byte-identical images for all 132 references (peak GPU memory
2.84\,GB). The IPAdapter-plus image embedding of each reference is precomputed on the CPU and
cached on disk keyed by the reference's SHA-256, so both ends read an identical tensor and the
2.5\,GB image encoder never has to share the GPU with Stable Diffusion.


\subsection{Phase 9: ControlNet zero-convolution fine-tuning}
Algorithm~2 of the paper fine-tunes only the zero-convolution layers of ControlNet. A short run
(50 steps, batch 1, \texttt{bfloat16}, calibration data only) was performed to validate the code
path; half precision produced NaN gradients within five steps, while \texttt{bfloat16} was stable.
A held-out probe loss improved from 0.1488 to 0.1296.

\subsection{Phase 10: EDICT}
EDICT was implemented from the coupled-transformation equations as its own module and tested in
isolation by inverting and regenerating random and image latents across a sweep of step counts
and mixing coefficients. At the operating point (20 steps, $p=0.93$,
half-precision UNet with single-precision latents) the maximum relative latent error was
$3.6\times10^{-3}$: a numerical floor from half precision, not an algorithmic failure. The sweep
also exposed a latent bug: the configuration included $p=0.5$, at which the coupling is
singular, so a full rebuild could never finish. The sweep now starts at 0.55.

\subsection{Phase 11: local hiding and recovery}
A full inversion treats the clean secret latent as pure noise and decodes to a reference-like
image (11.3\,dB against the secret), which is not imperceptible. Hiding therefore uses a
bounded-strength EDICT chain (strength 0.2, 20 steps, guidance 3.0) that is still an exact
paired inverse. Token noise and the sparse token mask are applied in the latent. IPAdapter
injects the reference embedding at both ends, making the stego depend on the reference pixels.

\subsection{Phase 12: AWGN channel}
The channel adds Gaussian noise whose strength is set from the signal power and the target
SNR, split equally across the in-phase and quadrature parts of complex symbols. A seeded generator
provides the noise realisation so that experiments are repeatable. The measured noise variance
matched the analytic value within 0.46\,\%.

\subsection{Phase 13: semantic codec}
The official SwinJSCC release with bandwidth $C=64$ trained at 10\,dB (the nearest public
checkpoint to the paper's $C=48$) was integrated in single precision, with only its noise
generator replaced by a seeded one. In isolation it reconstructs images at 32.54\,dB at 10\,dB
SNR, compared with 22.88\,dB for the VAE-based substitute used before the official weights were
obtained.

\subsection{Phase 14: CodeFormer enhancement}
The official CodeFormer architecture and checkpoint run at $512$\,px with fidelity 0.5 on the
facial class only, on the CPU because of GPU memory. The gate requires measured deltas rather
than assuming improvement; CodeFormer may lower PSNR slightly while improving LPIPS.

\subsection{Phase 15: assembled baseline}
The stages were assembled into one baseline object. A calibration run (15A) checks the regime
against the paper's Table~II with tolerances (Gate~A), and an evaluation run (15B) produces the
frozen table (Gate~B). The manifest records the configuration hash, split and token versions,
and the hash of every model. Results are in Chapter~6.

\subsection{Phase 16: security and steganalysis}
Recovery is measured under three receiver conditions: the \emph{correct} token, a valid
\emph{wrong} token, and \emph{no token} (a fixed public fallback token with an all-zero mask).
A XuNet steganalyser with a fixed SRM high-pass first layer is trained to separate clean images
from stego images, with the train, validation and test splits grouped by image so that patches of
one image never cross splits.

\section{RACS Reference Retrieval (Phases 17--22)}
All retrieval logic lives in one module shared by transmitter and receiver.

\subsection{Phase 17: exact retrieval}
The corpus embeddings are indexed with an exact inner-product index. A FAISS
\texttt{IndexFlatIP} is built and saved in its own process; within the pipeline, the same exact
inner product is computed with NumPy, because FAISS and PyTorch load conflicting OpenMP runtimes
in one process on Windows. The gate verifies that FAISS and the in-process search return identical
top-50 rows and scores, that each corpus image retrieves itself first with score 1.0, and that the
row order matches the index file.

\subsection{Phase 18: token selection}
The token is hashed with SHA-256 together with a fixed label, and the result modulo $K$
gives the position of the chosen image in the top-$K$ list. The gate verifies identical output in separate processes, that $K=1$ always selects index 0, that
over the evaluation tokens all ten buckets are used and the distribution is near uniform, and
that the language's built-in hash is never called.

\subsection{Phase 19: structural representation}
Each corpus image is converted to a dilated binary edge map at $128$\,px and cached. Query
structure is derived from the public $X_{\mathrm{feat}}$ (segment boundaries or keypoint
coverage). The cache is deterministic, and the gate requires at least three cases in which the
CLIP ranking and the structural ranking disagree; five were found.

\subsection{Phase 20: reranking}
Each candidate gets a combined score: a weight $\alpha$ times its CLIP similarity plus a weight
$\beta$ times its edge-map overlap with the public structure. A sweep on
calibration data only measured, for the selected top-10 pool, the mean semantic similarity to the
secret, the mean edge IoU, and the fraction of images whose pool changed relative to CLIP alone
(Figure~\ref{fig:sweep}). The objective weights semantic similarity and edge IoU equally.

\begin{figure}[H]
	\centering
	\begin{tikzpicture}
		\begin{axis}[viz, width=0.75\textwidth, height=5.4cm,
			xlabel={Semantic weight $\alpha$ \ (structural weight $\beta = 1-\alpha$)},
			ylabel={Calibration objective},
			xtick={0,0.25,0.5,0.75,1}, xmin=-0.08, xmax=1.08,
			ymin=0.762, ymax=0.7745,
			yticklabel style={/pgf/number format/fixed, /pgf/number format/precision=3}]
			\addplot[vizA, mark=*, mark size=2pt] coordinates {(0,0.7648) (0.25,0.7690) (0.5,0.7717) (0.75,0.7723) (1,0.7640)};
			\addplot[only marks, mark=*, mark size=3.6pt, vizB, forget plot] coordinates {(0.75,0.7723)};
			\node[font=\scriptsize, anchor=south, text=black] at (axis cs:0.75,0.7726) {chosen: $\alpha=0.75$};
		\end{axis}
	\end{tikzpicture}
	\caption{Rerank weight sweep on calibration data only. Mixing in some structure beats pure CLIP ranking; $\alpha=0.75$ scores best and changes the top-10 pool for two thirds of the images.}
	\label{fig:sweep}
\end{figure}

$\alpha = 0.75$, $\beta = 0.25$ maximises the objective; it keeps most of the semantic quality of
CLIP alone while raising structural agreement, and it changes the candidate pool for two thirds
of the images. These weights were frozen before any evaluation run.

\subsection{Phase 21: RACS end to end}
RACS-Base subclasses the baseline and overrides only reference acquisition. Retrieval for the
transmitter and for the receiver run in two separate processes before the diffusion stack loads
(CLIP text encoding crashed when loaded beside it). The gate requires
$X_{\mathrm{ref}}^{\mathrm{tx}} = X_{\mathrm{ref}}^{\mathrm{rx}}$ for every evaluation image, compared
by identifier and file hash. Both systems are then measured through the same object and the same
evaluation function, differing only in where the reference comes from.

\subsection{Phase 22: controlled comparison}
A comparison script reads both systems' measurements on the same cases and tokens, refuses
calibration data, and reports quality, security, latency and memory side by side (Chapter~6).

\section{RACS Exact-Recovery Layer}
The layer is a self-contained module with four functions: compaction,
carrier preparation, embedding and extraction. A container
format stores a magic value, the length and the format of the payload, and the HMAC tag.
Evaluation generates 66 paper-style carriers (EDICT at strength 0.5 with a RACS-retrieved
reference), embeds every evaluation secret with four progressively stronger variants, writes each
stego to a real PNG file and reads it back before extraction. Two classical detectors were
implemented: Westfeld's chi-square test and Sample Pair Analysis. For each detector the AUC is
computed between each stego and the same carrier without payload, so 0.5 means the detector cannot
tell them apart.

Figure~\ref{fig:payload} shows how one secret becomes the hidden bit stream. The token is
stretched once by PBKDF2 into three independent values: an encryption key, a MAC key and the seed
that orders the embedding positions. Because the positions are chosen from texture computed on the
top five bits of each value, and a $\pm1$ change never alters those bits, the receiver recomputes
exactly the same positions from the stego image alone.

\begin{figure}[H]
	\centering
	\begin{tikzpicture}[scale=0.8, transform shape]
		\hyphenpenalty=10000\exhyphenpenalty=10000
		% key schedule (row 1)
		\node[src, text width=2.0cm] (tok) at (0,0) {Token $\tau$};
		\node[newblk, text width=2.4cm] (kdf) at (3.0,0) {\textbf{PBKDF2}\\200{,}000 rounds};
		\node[draw, thick, minimum height=0.8cm, minimum width=2.3cm, fill=gray!8, font=\footnotesize, anchor=west] (ek) at (5.4,0) {encryption key};
		\node[draw, thick, minimum height=0.8cm, minimum width=1.6cm, fill=gray!8, font=\footnotesize, right=0pt of ek] (mk) {MAC key};
		\node[draw, thick, minimum height=0.8cm, minimum width=2.1cm, fill=gray!8, font=\footnotesize, right=0pt of mk] (ps) {position seed};
		% payload container (row 2)
		\node[src, text width=2.0cm] (sec) at (0,-2.0) {Secret file};
		\node[newblk, text width=2.4cm] (cmp) at (3.0,-2.0) {\textbf{Compact}\\lossless WebP if smaller};
		\node[draw, thick, minimum height=0.8cm, minimum width=1.0cm, fill=gray!8, font=\footnotesize, anchor=west] (hdr) at (5.4,-2.0) {RACS};
		\node[draw, thick, minimum height=0.8cm, minimum width=1.1cm, fill=gray!8, font=\footnotesize, right=0pt of hdr] (len) {length};
		\node[draw, thick, minimum height=0.8cm, minimum width=1.8cm, fill=green!12, font=\footnotesize, right=0pt of len] (body) {file bytes};
		\node[src, text width=2.0cm] (png) at (13.0,-2.0) {Stego PNG};
		% processing (row 3)
		\node[newblk, text width=2.4cm] (enc) at (6.6,-4.1) {\textbf{Encrypt}\\SHAKE-256 stream (encryption key)};
		\node[newblk, text width=2.4cm] (mac) at (9.8,-4.1) {\textbf{Seal}\\32-byte HMAC tag (MAC key)};
		\node[newblk, text width=2.6cm] (emb) at (13.0,-4.1) {\textbf{Embed}\\$\pm1$ in most textured 50\,\% (position seed)};
		\draw[arr] (tok) -- (kdf);
		\draw[arr] (kdf) -- (ek);
		\draw[arr] (sec) -- (cmp);
		\draw[arr] (cmp) -- (hdr);
		\draw[arr] (len.south) -- (len.south |- enc.north);
		\draw[arr] (enc) -- (mac);
		\draw[arr] (mac) -- (emb);
		\draw[arr] (emb) -- (png);
	\end{tikzpicture}
	\caption{How the exact-recovery layer packs a secret. The token is stretched into three independent values, each used by one step; a wrong token produces a different MAC key, so the seal check fails and the file is rejected.}
	\label{fig:payload}
\end{figure}

\section{Prototype}
A web server on the laptop serves a mobile page. The sender uploads a photograph and receives a
stego PNG and an 8-digit key; the receiver uploads the PNG and the key and receives the secret.
Photographs keep their aspect ratio by edge-padding to a square at $512$\,px and cropping back on
receipt. The public-key generation and retrieval run in a separate worker process for the reasons
given above. A self-test mode runs a round trip from the command line. Figure~\ref{fig:sequence}
shows one complete exchange.

\begin{figure}[H]
	\centering
	\begin{tikzpicture}[scale=0.85, transform shape]
		\hyphenpenalty=10000\exhyphenpenalty=10000
		% participants
		\node[src, text width=2.4cm] (s) at (0,0) {\textbf{Sender phone}};
		\node[blk, text width=2.6cm] (l) at (5.5,0) {\textbf{Laptop server}};
		\node[src, text width=2.4cm] (r) at (11,0) {\textbf{Receiver phone}};
		\foreach \n in {s,l,r} \draw[dashed, gray!70] (\n.south) -- ++(0,-9.2);
		% send
		\draw[arr] (0,-1.1) -- node[above, font=\footnotesize]{upload secret photo} (5.5,-1.1);
		\node[newblk, text width=3.6cm, font=\footnotesize] at (5.5,-2.1) {$K_{\mathrm{pub}}$, RACS retrieval, EDICT stego look, encrypt + embed};
		\draw[arr] (5.5,-3.5) -- node[above, font=\footnotesize]{stego PNG + 8-digit key} (0,-3.5);
		% hand-over between people
		\draw[arr] (0,-4.4) -- node[above, font=\footnotesize]{PNG sent as a file (any messenger)} (11,-4.4);
		\draw[arr, dashed] (0,-5.1) -- node[above, font=\footnotesize]{key shared separately} (11,-5.1);
		% receive
		\draw[arr] (11,-6.0) -- node[above, font=\footnotesize]{upload PNG + key} (5.5,-6.0);
		\node[newblk, text width=3.6cm, font=\footnotesize] at (5.5,-7.0) {read $\pm1$ bits, check seal, decrypt (no AI model)};
		\draw[arr] (5.5,-8.2) -- node[above, font=\footnotesize]{exact secret, or ``rejected''} (11,-8.2);
	\end{tikzpicture}
	\caption{One exchange in the phone prototype. The laptop does all computation; a recompressed photo or a wrong key is detected by the seal and rejected.}
	\label{fig:sequence}
\end{figure}

\section{Verification Results}
Every phase writes its gate result to a JSON record of named checks, so the state of the whole
build can be audited at any time. Figure~\ref{fig:gates} summarises the stored records: 119 checks
across twelve phases, 113 of which pass. All six failures sit in Phases 15 and 16 and have a
single cause: those phases were run on the 2-image-per-class subset, while their gates expect the
full 22 images per class (sample counts, record counts, and the hash of the full frozen table).
Every determinism, provenance and component check passes, including all RACS phases.

\begin{figure}[H]
	\centering
	\begin{tikzpicture}
		\begin{axis}[vizbar, ybar stacked, bar width=13pt, width=0.95\textwidth, height=5.6cm,
			symbolic x coords={5,7,8,9,10,15,16,17,18,19,20,21}, xtick=data,
			xlabel={Phase}, ylabel={Verification checks}, ymin=0, ymax=22,
			enlarge x limits=0.05,
			legend style={at={(0.98,0.97)}, anchor=north east}]
			\addplot[fill=green!55!black, draw=none] coordinates {(5,20) (7,13) (8,12) (9,10) (10,14) (15,11) (16,12) (17,6) (18,5) (19,3) (20,3) (21,4)};
			\addplot[fill=red!70!black, draw=none] coordinates {(5,0) (7,0) (8,0) (9,0) (10,0) (15,1) (16,5) (17,0) (18,0) (19,0) (20,0) (21,0)};
			\legend{Passed, Failed (subset size only)}
		\end{axis}
	\end{tikzpicture}
	\caption{Automatic verification checks per phase, from the stored gate records. Phases 17--21 are the RACS phases.}
	\label{fig:gates}
\end{figure}

\section{Engineering Challenges}
\begin{itemize}
	\item \textbf{Windows commit charge.} Memory-mapped checkpoint loading reserved more virtual
	memory than the paging file allowed, causing access violations. A small shim loads
	\texttt{safetensors} files read-only before any model library is imported.
	\item \textbf{6\,GB GPU.} Half precision, attention slicing, $256$\,px operation, one model
	resident at a time, CPU placement of the IPAdapter encoder and CodeFormer, and on-disk caches
	of derived tensors kept every phase within budget.
	\item \textbf{Cross-process determinism.} Process-randomised hashing, half-precision encoders
	and implicit device nondeterminism were removed from every path that affects the reference.
	\item \textbf{Float-to-integer truncation.} A conversion on the transmitter side truncated
	instead of rounding before the IPAdapter embedding, which could have desynchronised the two
	ends; caching by reference hash removed the risk.
	\item \textbf{Mixed-state detection.} Upgrading components invalidated earlier frozen
	artifacts. Configuration and model hashes in every manifest made the verification gates fail
	loudly until the affected phases were rebuilt, instead of silently mixing results.
\end{itemize}
